\section*{Chapitre \ref{ch:g7:shadows-eclipses} --- Ombres et éclipses}

\begin{solution}{exo:g7:shadows-eclipses:1}
\omterm{prop:g7:shadows-eclipses:umbra}{Zone d'ombre}~: la source est entièrement cachée --- aucune \omterm{def:g1:light-and-shadows:source}{lumière}
d'aucune de ses parties n'arrive. \omterm{prop:g7:shadows-eclipses:umbra}{Zone de pénombre}~: la source est
partiellement cachée --- la \omterm{def:g1:light-and-shadows:source}{lumière} d'une partie d'elle arrive
encore, en dépassant par-dessus le bord de l'obstacle.
\end{solution}

\begin{solution}{exo:g7:shadows-eclipses:2}
Un point n'a pas de «~parties~» que l'on puisse cacher à moitié~:
chaque endroit le voit ou ne le voit pas --- \omterm{def:g1:light-and-shadows:shadow}{ombre} nette, aucune
frange. Un ciel couvert est une source si étendue que presque chaque
endroit en voit la plus grande partie~: quasiment toute pénombre, et
des \omterm{def:g1:light-and-shadows:shadow}{ombres} effacées.
\end{solution}

\begin{solution}{exo:g7:shadows-eclipses:3}
Les rayons issus des \emph{bords} de la source, qui frôlent les côtés
de l'obstacle. La \omterm{prop:g7:shadows-eclipses:umbra}{zone d'ombre} est la région que les rayons des deux
bords manquent.
\end{solution}

\begin{solution}{exo:g7:shadows-eclipses:4}
De Soleil~: Soleil--Lune--Terre, dans cet ordre --- la \omterm{def:g2:moon-in-the-sky:moon}{Lune} doit être
nouvelle (entre le Soleil et nous). De \omterm{def:g2:moon-in-the-sky:moon}{Lune}~: Soleil--Terre--Lune ---
la \omterm{def:g2:moon-in-the-sky:moon}{Lune} doit être pleine (à l'opposé du Soleil).
\end{solution}

\begin{solution}{exo:g7:shadows-eclipses:5}
La totalité~: seulement ceux qui se trouvent dans la petite tache
filante de la \omterm{prop:g7:shadows-eclipses:umbra}{zone d'ombre}. Une morsure partielle~: le vaste
pourtour de pénombre. Absolument rien~: tous ceux que les régions
d'\omterm{def:g1:light-and-shadows:shadow}{ombre} manquent --- la plus grande partie de la \omterm{def:g4:solar-system:planet}{planète}.
\end{solution}

\begin{solution}{exo:g7:shadows-eclipses:6}
L'\omterm{def:g7:shadows-eclipses:lunar}{éclipse de Lune} est la \omterm{def:g2:moon-in-the-sky:moon}{Lune} elle-même qui s'assombrit --- un
changement sur un astre que toute la face nocturne peut voir d'un
coup. La totalité solaire exige de se tenir \emph{à l'intérieur} de
la mince \omterm{prop:g7:shadows-eclipses:umbra}{zone d'ombre} de la \omterm{def:g2:moon-in-the-sky:moon}{Lune}, là où elle touche le sol~: une
tache, non un événement à l'échelle du ciel.
\end{solution}

\begin{solution}{exo:g7:shadows-eclipses:7}
L'\omterm{def:g6:sun-earth-moon:axisorbit}{orbite} de la \omterm{def:g2:moon-in-the-sky:moon}{Lune} est inclinée d'environ cinq degrés sur le plan
Terre--Soleil~: la plupart des nouvelles et des \omterm{ex:g2:moon-in-the-sky:shapes}{pleines lunes} passent
au-dessus ou au-dessous de la ligne exacte, et les \omterm{def:g1:light-and-shadows:shadow}{ombres} manquent
leur cible. Ce n'est que lorsque la phase et la traversée du plan
coïncident que les trois mondes s'alignent.
\end{solution}

\begin{solution}{exo:g7:shadows-eclipses:8}
La projection par sténopé sur du papier, et les croissants au sol
sous un arbre feuillu (ou une passoire)~; des lunettes d'éclipse
certifiées et intactes pour l'observation directe. L'œil nu~:
uniquement pendant la totalité elle-même, et le premier éclat qui
revient met fin à la permission.
\end{solution}

\begin{solution}{exo:g7:shadows-eclipses:9}
Le Soleil fait environ $400$ fois la largeur de la \omterm{def:g2:moon-in-the-sky:moon}{Lune}, à environ
$400$ fois sa distance~: les deux disques se correspondent donc dans
le ciel presque exactement. Enfants de cet ajustement~: la totalité
existe, mais tout juste --- brève, avec une minuscule tache
d'\omterm{def:g1:light-and-shadows:shadow}{ombre}~; et quand la \omterm{def:g2:moon-in-the-sky:moon}{Lune} parcourt la partie éloignée de son \omterm{def:g6:sun-earth-moon:axisorbit}{orbite}
ovale, son disque légèrement plus petit laisse le bord du Soleil
découvert --- c'est l'anneau de l'éclipse annulaire.
\end{solution}

\begin{solution}{exo:g7:shadows-eclipses:10}
L'ourlet de pénombre de toute \omterm{def:g1:light-and-shadows:shadow}{ombre} est tracé par les bords du
Soleil, et la construction montre que cet ourlet s'élargit avec
l'écart entre l'obstacle et l'écran~: les rayons issus de bords
solaires opposés s'écartent d'autant plus qu'ils voyagent longtemps
après l'obstacle. Les pieds sont à écart nul --- bord net~; l'\omterm{def:g1:light-and-shadows:shadow}{ombre}
de la tête se trouve deux \omterm{def:g2:measuring-length:units}{mètres} d'écartement plus loin --- bord
doux.
\end{solution}

\begin{solution}{exo:g7:shadows-eclipses:11}
Chaque trouée entre les feuilles est une chambre noire à sténopé
naturelle~; chaque tache au sol est l'\omterm{prop:g5:mirrors-reflection:image}{image} du Soleil formée par ce
sténopé --- ronde les jours ordinaires. Pendant l'éclipse partielle,
le Soleil lui-même est un croissant~: chaque petite chambre fidèle
projette donc un croissant --- et tous pointent dans le même sens
parce que tous copient l'unique croissant du ciel.
\end{solution}

\begin{solution}{exo:g7:shadows-eclipses:12}
La \omterm{def:g1:light-and-shadows:source}{lumière} du Soleil qui frôle le limbe de la Terre traverse
l'\omterm{def:g5:air-pressure-weather:pressure}{atmosphère} par ses plus longues routes possibles --- des trajets de
longueur de coucher de Soleil, tout autour de la \omterm{def:g4:solar-system:planet}{planète} à la fois.
L'\omterm{def:g2:air-around-us:air}{air} diffuse une grande part de la \omterm{def:g1:light-and-shadows:source}{lumière} qui passe (celle-là même
qui peint les ciels de la Terre), et la \omterm{def:g7:light-sources-propagation:diffusion}{diffusion} dérobe en premier
les couleurs de court trajet~: la \omterm{def:g1:light-and-shadows:source}{lumière} qui survit au long voyage
rasant et poursuit jusqu'à la \omterm{def:g2:moon-in-the-sky:moon}{Lune} est donc la survivante des
couchers de Soleil --- rouge. La \omterm{def:g2:moon-in-the-sky:moon}{Lune} éclipsée luit de la \omterm{def:g1:light-and-shadows:source}{lumière} de
tous les levers et de tous les couchers de Soleil de la Terre tombant
ensemble sur elle.
\end{solution}

\begin{solution}{pb:g7:shadows-eclipses:1}
\textbf{1.} La bande est le trajet de la \omterm{prop:g7:shadows-eclipses:umbra}{zone d'ombre} de la \omterm{def:g2:moon-in-the-sky:moon}{Lune} qui
file sur le sol~; la large zone qui la flanque est l'empreinte de la
pénombre.
\textbf{2.} Une éclipse partielle~: la ville se trouve dans la
pénombre --- une morsure dans le Soleil, observée avec une
protection convenable, et aucune obscurité.
\textbf{3.} Même un croissant de Soleil de cinq pour cent est des
milliers de fois plus brillant que la \omterm{ex:g2:moon-in-the-sky:shapes}{pleine lune}~: il maintient le
ciel éclairé et la couronne invisible. La nuit à midi appartient à la
seule \omterm{prop:g7:shadows-eclipses:umbra}{zone d'ombre} --- il n'y a pas de «~presque totalité~».
\textbf{4.} La tache d'\omterm{def:g1:light-and-shadows:shadow}{ombre} est à peu près ronde~: un campement sur
l'\omterm{def:g6:sun-earth-moon:axisorbit}{axe} se trouve sous toute sa largeur pendant qu'elle file, alors que
les campements proches du bord de la bande n'en attrapent qu'une
corde du cercle --- une traversée plus courte, une totalité plus
courte.
\textbf{5.} Les lunettes d'éclipse certifiées bloquent la brûlure du
Soleil~; les lunettes de soleil, non --- elles atténuent
l'éblouissement tout en laissant passer la part qui brûle la rétine,
ce qui rend possible une contemplation confortable~: pire que rien.
\textbf{6.} Un sténopé à une extrémité, un écran de papier calque à
l'autre (ou du papier ordinaire posé au sol)~: se tenir \emph{dos} au
Soleil, le trou pointé par-dessus l'épaule, et regarder le disque
projeté --- jamais à travers le trou.
\textbf{7.} Rejetée~: le reflet de l'eau est une \omterm{prop:g5:mirrors-reflection:image}{image} du Soleil dans
un \omterm{def:g5:mirrors-reflection:reflection}{miroir}, presque aussi aveuglante que le Soleil lui-même --- une
brûlure rebondie reste une brûlure.
\textbf{8.} Des dizaines d'\omterm{prop:g5:mirrors-reflection:image}{images} en croissant sur le drap~: chaque
trou est une chambre à sténopé, et chacune projette le Soleil mordu
--- une passoire pleine d'éclipses.
\textbf{9.} La \omterm{def:g2:moon-in-the-sky:moon}{Lune} --- invisible dans l'éblouissement jusqu'à ce que
le bord de son disque commence à recouvrir celui du Soleil~; le
campement est entré dans la pénombre.
\textbf{10.} Le Soleil visible se réduit à un mince croissant --- une
source de plus en plus proche du point. Source plus petite, ourlets
de pénombre plus étroits~: le bord doux de chaque \omterm{def:g1:light-and-shadows:shadow}{ombre} se resserre,
et le monde devient étrange et net.
\textbf{11.} La \omterm{prop:g7:shadows-eclipses:umbra}{zone d'ombre} de la \omterm{def:g2:moon-in-the-sky:moon}{Lune} elle-même --- la classe se
tient à l'intérieur de la pointe du cône d'\omterm{def:g1:light-and-shadows:shadow}{ombre} pendant qu'elle file
au-dessus du campement. Autour du disque noir apparaît la couronne,
l'\omterm{def:g5:air-pressure-weather:pressure}{atmosphère} externe et nacrée du Soleil, visible seulement
maintenant~; le premier éclat du disque qui revient rétablit les
règles du plein jour --- lunettes remises.
\textbf{12.} Par exemple~: «~Une éclipse est une \omterm{def:g1:light-and-shadows:shadow}{ombre} dont la lampe,
l'obstacle et l'écran sont des mondes~: la \omterm{prop:g7:shadows-eclipses:umbra}{zone d'ombre} de la \omterm{def:g2:moon-in-the-sky:moon}{Lune}
qui effleure notre sol, ou la \omterm{def:g2:moon-in-the-sky:moon}{Lune} qui navigue à travers la nôtre ---
une géométrie que l'on peut dessiner avec deux rayons et une règle,
mise en scène à la plus grande échelle de la nature.~»
\end{solution}
